\chapter{Estado del arte}
\label{chap:estado_arte}


\section{Ingeniería de Líneas de Productos Software (SPL)}
\label{sec:ingenieria_spl}

\subsection{Líneas de Productos Software}
\label{subsec:spl}

A lo largo de la historia, el desarrollo de software se ha basado tradicionalmente en la construcción de sistemas individuales orientados a resolver necesidades concretas, donde cada nuevo producto software suele desarrollarse como un proyecto independiente, incluso cuando comparte una parte significativa de sus funcionalidades, arquitectura o componentes con otros sistemas existentes. Esta estrategia puede resultar adecuada cuando el número de productos es reducido; sin embargo, hoy en día esto presenta limitaciones importantes para la gran mayoría de organizaciones, ya que surge de forma paralela la necesidad de adaptar un determinado sistema base a una gran cantidad de clientes con problemas ligera o enormemente diferentes.

Debido a este hecho surge la Ingeniería de Líneas de Productos Software. Clements y Northrop definen las SPL de la siguiente forma: \enquote{\textit{a set of software-intensive systems that share a common, managed set of features satisfying the specific needs of a particular market segment or mission and that are developed from a common set of core assets in a prescribed way}} \cite{clements2002}. Gracias a ellas es posible identificar, gestionar y reutilizar aquellos elementos comunes entre diferentes productos, a la vez que se mantienen explícitamente aquellos aspectos que pueden variar entre las distintas configuraciones posibles, haciendo así factible y mucho más viable la gestión de dichos elementos en lugar de construir cada producto desde cero.

La Figura~\ref{fig:ejemplo_spl} muestra una forma de entender este concepto.

\begin{figure}[H]
    \centering
    \includegraphics[width=0.6\textwidth]{SPL.jpg}
    \caption{Representación de una línea de productos software basada en activos comunes. Extraído de \cite{clements2002}.}
    \label{fig:ejemplo_spl}
\end{figure}

La existencia de elementos comunes y diferencias entre los productos introduce el concepto de variabilidad, definido por van Gurp, Bosch y Svahnberg como \enquote{\textit{Variability is the ability to change or customize a system}} \cite{vanGurp2001}. Esta debe ser identificada, representada y gestionada durante el proceso de desarrollo, ya que permite definir qué elementos pueden formar parte de cada producto y bajo qué condiciones pueden combinarse.

La gestión de esta variabilidad requiere distinguir entre dos actividades complementarias. Por una parte, la ingeniería de dominio (\textit{domain engineering}) se encarga de analizar un determinado dominio, identificar sus elementos comunes y variables, y construir los activos reutilizables que formarán la base de la línea de productos. Por otra parte, la ingeniería de aplicación (\textit{application engineering}) utiliza dichos activos para derivar productos concretos mediante la selección de las características necesarias y la resolución de las variabilidades definidas.

No obstante, este paradigma introduce consigo retos relacionados con la representación y gestión de la complejidad. A medida que aumenta el número de características variables y configuraciones posibles, resulta necesario disponer de mecanismos que permitan gestionar esta variabilidad de forma íntegra y precisa, lo cual da surgimiento a técnicas como los modelos de características.


\subsection{Modelos de características}
\label{subsec:feature_models}

Los modelos de características (\textit{Feature Models}) son considerados uno de los mecanismos más utilizados para representar la variabilidad dentro de las líneas de productos software. Originalmente, fueron introducidos como parte del método FODA (\textit{Feature-Oriented Domain Analysis}) por Kang et al., con el objetivo de describir las características comunes y variables existentes dentro de una familia de sistemas software \cite{kang1990}. Benavides et al. los describen de la siguiente forma: \enquote{\textit{a feature model represents the information of all possible products of a software product line in terms of features and relationships among them}} \cite{benavides2010}. Estos modelos permiten representar el espacio de configuraciones de una línea de productos mediante las características que pueden formar parte de cada producto y las relaciones existentes entre ellas.

Para entenderlo mejor, estos modelos se pueden representar mediante una estructura jerárquica en forma de árbol, donde cada nodo representa una característica del sistema. Las relaciones existentes entre características permiten establecer dependencias entre ellas, indicando qué elementos son obligatorios, opcionales o alternativas dentro de una configuración concreta. Además, pueden incorporar restricciones adicionales que expresan condiciones que no pueden representarse únicamente mediante la jerarquía del modelo. Véase la Figura~\ref{fig:ejemplo_feature_model}.

\begin{figure}[H]
    \centering
    \includegraphics[width=1\textwidth]{ejemplo_feature_model.png}
    \caption{Ejemplo de modelo de características. Extraído de \cite{uvl}.}
    \label{fig:ejemplo_feature_model}
\end{figure}

Aunque los modelos de características tradicionales permiten representar una gran cantidad de escenarios de variabilidad, inicialmente estaban orientados principalmente a expresar relaciones estructurales entre características. Con el aumento de la complejidad de los sistemas software surgió la necesidad de incorporar mecanismos adicionales, como atributos, restricciones aritméticas o diferentes tipos de dominios, que permitan representar escenarios más cercanos a los sistemas reales. A su vez, al incrementar la expresividad de los modelos, también se incrementa la complejidad de los mismos, lo cual hace necesaria la existencia de herramientas capaces de procesar estos modelos automáticamente y verificar propiedades sobre las configuraciones posibles.

\subsection{Análisis automático y motores de razonamiento}
\label{subsec:analisis_automatico_motores_razonamiento}

Una vez definidos los modelos de características como mecanismo para representar la variabilidad de una SPL, surge la necesidad de disponer de técnicas capaces de analizar automáticamente la información contenida en dichos modelos. El análisis automático permite obtener información sobre las configuraciones representadas, comprobar propiedades del modelo y realizar operaciones de razonamiento sin necesidad de explorar manualmente cada una de las posibles variantes.

Los motores de razonamiento (\textit{reasoning engines}) se encargan de
analizar automáticamente los modelos de variabilidad. Entre otras operaciones,
permiten comprobar si un modelo es satisfacible, detectar inconsistencias e
identificar características que no pueden formar parte de ninguna configuración
válida.

Por ejemplo, puede considerarse un modelo de características cuya raíz es
\textit{Primer Plato}. Las características \textit{Pasta} y \textit{Arroz}
forman una relación alternativa. Esto significa que, en cada configuración
válida, debe seleccionarse exactamente una de las características que forman
parte de dicho grupo: puede elegirse \textit{Pasta} o \textit{Arroz}, pero no
ambas a la vez. Por su parte, \textit{Salsa} se define como una característica
opcional de \textit{Primer Plato}, por lo que puede seleccionarse o no junto
con la alternativa elegida. La Figura~\ref{fig:ejemplo_sat} muestra esta
estructura.

\begin{figure}[H]
    \centering
    \includegraphics[width=0.6\textwidth]{ejemplo_fm_propio.png}
    \caption{Ejemplo gráfico de modelo de características. Elaboración propia.}
    \label{fig:ejemplo_sat}
\end{figure}

Si se añade como restricción global la expresión \texttt{Pasta \& Arroz}, el
modelo resulta insatisfacible, ya que la restricción exige seleccionar
simultáneamente dos características que la relación alternativa impide
seleccionar a la vez. La expresión puede escribirse como una restricción
válida; sin embargo, no existe ninguna configuración que satisfaga al mismo
tiempo la jerarquía del modelo y dicha restricción.

Una de las aproximaciones más utilizadas para realizar análisis automático sobre modelos de características consiste en transformar dichos modelos en fórmulas proposicionales que puedan ser procesadas mediante solucionadores SAT (\textit{Boolean Satisfiability Problem}). El problema SAT consiste en determinar si existe una asignación de valores para un conjunto de variables booleanas que haga verdadera una fórmula lógica dada. En el contexto de las líneas de productos software, las características y restricciones de un modelo pueden representarse mediante variables booleanas y expresiones lógicas equivalentes, permitiendo utilizar solucionadores SAT para responder consultas relacionadas con la validez de configuraciones o la existencia de productos derivados \cite{benavides2010}.

Para aplicar esta técnica, el modelo se traduce normalmente a una fórmula lógica equivalente que conserva las restricciones definidas entre sus elementos. Por ejemplo, una relación obligatoria entre dos características puede expresarse mediante una implicación lógica, mientras que las restricciones de exclusión entre características pueden representarse mediante cláusulas que impidan determinadas combinaciones. Una vez realizada esta transformación, el solucionador SAT puede determinar si existe una configuración válida o identificar aquellas asignaciones que cumplen las restricciones del modelo.

Además de la codificación SAT, los Diagramas de Decisión Binarios
(\textit{Binary Decision Diagrams}, BDD) permiten representar una función
booleana mediante un grafo dirigido acíclico. Cada nodo interno corresponde a
una variable de decisión y cada nodo terminal representa uno de los valores de
la función. Esta representación permite codificar de forma compacta el espacio
de configuraciones válidas y realizar distintas operaciones de análisis sobre
él \cite{bryant1986}.

Para ilustrar ambas representaciones sobre un mismo ejemplo, considérese la
fórmula $\varphi = (A \Rightarrow B) \land \neg(A \land C)$. La primera
restricción establece que la selección de $A$ requiere la de $B$, mientras que
la segunda impide seleccionar simultáneamente $A$ y $C$. Su forma normal
conjuntiva es
$\varphi \equiv (\neg A \lor B) \land (\neg A \lor \neg C)$.

La Figura~\ref{fig:sat_vs_bdd} muestra la representación de esta misma fórmula
mediante su CNF (Forma Normal Conjuntiva) para su análisis con SAT y mediante un
BDD reducido. En el BDD, cada arista etiquetada con $0$ o $1$ representa la
asignación de los valores falso y verdadero, respectivamente, a la variable del nodo de
origen, mientras que los terminales $0$ y $1$ indican que la fórmula resulta
falsa o verdadera. De esta forma, ambas partes de la figura utilizan
exactamente las mismas restricciones y variables.

\begin{figure}[H]
    \centering
    \makebox[\textwidth][c]{%
        \includegraphics[width=1.1\textwidth]{sat_vs_bdd.png}
    }
    \caption{Representación de una misma fórmula mediante CNF y un BDD reducido. Elaboración propia.}
    \label{fig:sat_vs_bdd}
\end{figure}

La Tabla~\ref{tab:sat_vs_bdd} compara SAT y BDD como técnicas de análisis automático. SAT permite resolver consultas de satisfacibilidad, mientras que un BDD representa la función booleana del modelo y puede reutilizarse para distintas operaciones.

\begin{table}[H]
    \centering
    \begin{tabular}{>{\raggedright\arraybackslash}p{4cm}
                    >{\raggedright\arraybackslash}p{5cm}
                    >{\raggedright\arraybackslash}p{5cm}}
        \toprule
        \textbf{Criterio} & \textbf{SAT} & \textbf{BDD} \\
        \midrule
        Representación del problema &
        Transformación del modelo de características a una fórmula lógica que debe ser evaluada por un solucionador SAT. &
        Construcción de una estructura gráfica que representa las posibles asignaciones de la función booleana. \\
        \midrule
        Tipo de análisis &
        Orientado principalmente a responder consultas concretas sobre la existencia de configuraciones válidas. &
        Permite realizar análisis globales sobre el espacio completo de configuraciones representado. \\
        \midrule
        Resultado obtenido &
        Determina si existe una asignación que satisface las restricciones y, en algunos casos, proporciona una configuración válida. &
        Proporciona una representación compacta del conjunto de configuraciones válidas. \\
        \midrule
        Ventaja principal &
        Gran eficiencia para comprobar propiedades concretas del modelo y resolver problemas de satisfacibilidad. &
        Permite realizar múltiples operaciones de análisis sobre el modelo una vez construido el BDD. \\
        \midrule
        Limitación principal &
        Requiere resolver consultas individuales y no mantiene necesariamente una representación completa del espacio de configuraciones. &
        La construcción del diagrama puede resultar costosa para modelos con elevada complejidad. \\
        \bottomrule
    \end{tabular}
    \caption{Comparativa entre SAT y BDD como técnicas de análisis automático de modelos de características.}
    \label{tab:sat_vs_bdd}
\end{table}


\section{Universal Variability Language (UVL)}
\label{sec:uvl}

Universal Variability Language (UVL) es un lenguaje textual para representar
modelos de características o \textit{feature models}, así como sus
restricciones de variabilidad \cite{uvl,uvlparser2023}. En este trabajo se
utiliza como formato de salida del generador, por lo que esta sección se centra
en las construcciones que el sistema puede producir y que posteriormente son
procesadas por las herramientas del ecosistema Flamapy.

El Apéndice~\ref{app:modelo-uvl} incluye ejemplos más extensos de modelo UVL
completo. Los modelos presentados en esta sección también se han construido
como archivos UVL completos y se han validado mediante el \textit{parser}
utilizado por Flamapy.

\subsection{Estructura básica de un modelo UVL}
\label{subsec:estructura_basica}

Un archivo UVL generado por el sistema se compone principalmente de los
bloques \texttt{include}, \texttt{features} y \texttt{constraints}. El bloque
\texttt{include}, situado al comienzo del archivo cuando es necesario, declara
los \textit{minor levels} de UVL utilizados por el modelo. El bloque
\texttt{features} define el árbol de características o \textit{feature tree},
mientras que \texttt{constraints} recoge las restricciones que relacionan
elementos de distintas ramas de dicho árbol.

Cada modelo tiene una única característica raíz o \textit{root feature}. A
partir de ella, las características se organizan mediante indentación y se
agrupan según la relación de selección que exista entre una característica
padre y sus características hijas. Los nombres que contienen espacios o
caracteres especiales se escriben entre comillas dobles, y no deben confundirse con valores de tipo cadena de atributos o características de este tipo, que se escriben con comillas simples.

\subsection{Relaciones y \textit{Boolean level}}
\label{subsec:diseño_lenguaje}

El \textit{Boolean level} constituye la base de los modelos generados. Permite
definir las relaciones estructurales de un \textit{feature tree} mediante los
grupos \textit{mandatory}, \textit{optional}, \textit{alternative} y
\textit{or}.

Una característica \textit{mandatory} debe seleccionarse siempre que se haya
seleccionado su característica padre. En cambio, una característica
\textit{optional} puede seleccionarse o no cuando la característica padre
forma parte de la configuración. Un grupo \textit{alternative} exige
seleccionar exactamente una de sus características hijas, mientras que un
grupo \textit{or} exige seleccionar al menos una de ellas y permite seleccionar
varias simultáneamente.

Además de estas relaciones jerárquicas, UVL permite definir restricciones entre
características pertenecientes a partes distintas del árbol. Estas restricciones
se denominan \textit{cross-tree constraints} y pueden emplear operadores de
negación (\texttt{!}), conjunción (\texttt{\&}), disyunción (\texttt{|}),
implicación (\texttt{=>}) y equivalencia (\texttt{<=>}).

El generador también puede activar el \textit{minor level} de
\textit{group cardinality}. Esta construcción limita el número de
características hijas seleccionables dentro de un grupo mediante una expresión
de la forma \texttt{[mínimo..máximo]}. Por ejemplo, una cardinalidad
\texttt{[0..2]} permite seleccionar desde ninguna hasta dos características
del grupo.

El Listado~\ref{lst:modelo-uvl-booleano} muestra un modelo completo del
\textit{Boolean level}. La característica \texttt{Accessories} utiliza una
\textit{group cardinality} de \texttt{[0..2]}, por lo que pueden seleccionarse
como máximo dos correas. La primera restricción entre árboles (\textit{cross-tree constraint}) exige que
el modelo incluya conexión mediante Wi-Fi o Bluetooth, mientras que la segunda
establece que la resistencia al agua requiere disponer de pantalla.

\begin{lstlisting}[
    language={},
    caption={Modelo UVL completo del \textit{Boolean level} con \textit{group cardinality}.},
    label={lst:modelo-uvl-booleano}
]
include
    Boolean.group-cardinality

features
    Smartwatch
        optional
            GPS
            Display
            "Wi-Fi"
            Bluetooth
            "Water Resistance"
            Accessories
                [0..2]
                    "Leather Strap"
                    "Metal Strap"
                    "Silicone Strap"

constraints
    "Wi-Fi" | Bluetooth
    "Water Resistance" => Display
\end{lstlisting}

\subsection{Niveles de expresividad avanzados}
\label{subsec:niveles_lenguaje}

Los niveles de UVL permiten incorporar construcciones adicionales sobre el
núcleo booleano. El \textit{Arithmetic level} añade \textit{attributes}
numéricos y operaciones sobre sus valores, mientras que el \textit{Type level}
permite declarar características con valores de tipo \texttt{Integer},
\texttt{Real} o \texttt{String}. Cada nivel incluye las capacidades de los
niveles anteriores.

\subsubsection{\textit{Arithmetic level}, \textit{attributes} y niveles menores}

Un \textit{attribute} asocia un valor a una característica. En UVL se declara
entre llaves después del nombre de dicha característica. Por ejemplo,
\texttt{Standard \{Capacity 300\}} asigna el valor 300 al \textit{attribute}
\texttt{Capacity} de la característica \texttt{Standard}. Estos valores pueden
utilizarse posteriormente en las restricciones del modelo.

El \textit{Arithmetic level} permite aplicar operadores como \texttt{+},
\texttt{-}, \texttt{*}, \texttt{/}, \texttt{==}, \texttt{!=}, \texttt{<} o
\texttt{>=} sobre \textit{attributes} numéricos. El generador puede incorporar
además los \textit{minor levels} de \textit{feature cardinality} y
\textit{aggregate functions}.

La \textit{feature cardinality}, declarada mediante
\texttt{Arithmetic.feature-cardinality}, permite indicar cuántas instancias de
una característica pueden seleccionarse. Por su parte, las
\textit{aggregate functions} permiten operar sobre los valores de un mismo
\textit{attribute} presentes en las características seleccionadas. La función
\texttt{sum()} calcula la suma de dichos valores y \texttt{avg()} calcula su
media; por tanto, no cuentan características, sino que agregan los valores de
sus \textit{attributes}.

El Listado~\ref{lst:modelo-uvl-aritmetico} combina estas construcciones en un
modelo completo. \texttt{Storage} es un \textit{attribute} asociado a los
reproductores multimedia, mientras que \texttt{Capacity} se asocia a las
opciones de batería. La primera restricción exige que la suma de los valores de
almacenamiento sea 12 y la segunda establece que la capacidad media de las
baterías disponibles debe ser, como mínimo, 300.

\begin{lstlisting}[
    language={},
    caption={Modelo UVL completo con \textit{attributes}, \textit{feature cardinality} y \textit{aggregate functions}.},
    label={lst:modelo-uvl-aritmetico}
]
include
    Arithmetic.feature-cardinality
    Arithmetic.aggregate-function

features
    Smartwatch
        mandatory
            Connectivity
                or
                    Bluetooth
                    "Wi-Fi"
            Battery
                alternative
                    Standard {Capacity 300}
                    "Long Life" {Capacity 450}
        optional
            "Music Player" {Storage 2}
            "Film Player" {Storage 10}
            "Watch Faces" cardinality [5..20]

constraints
    ("Music Player".Storage + "Film Player".Storage) == 12
    avg(Capacity) >= 300
\end{lstlisting}

\subsubsection{\textit{Type level} y \textit{string constraints}}

El \textit{Type level} amplía el modelo mediante características cuyo valor
pertenece a un tipo determinado. UVL permite declarar características de tipo
\texttt{Integer}, \texttt{Real} y \texttt{String}; el tipo se escribe antes
del nombre de la característica. Estas declaraciones no deben confundirse con
los \textit{attributes}: una característica tipada representa directamente un
valor configurable, mientras que un \textit{attribute} añade información a una
característica del modelo.

El \textit{minor level} de \textit{string constraints}, activado mediante
\texttt{Type.string-constraints}, permite expresar restricciones sobre valores
de tipo cadena. En particular, la función \texttt{len()} obtiene la longitud
de una cadena, por lo que puede utilizarse para imponer un número concreto de
caracteres o establecer límites sobre dicho valor.

El Listado~\ref{lst:modelo-uvl-tipos} muestra un modelo completo con
características de tipo \texttt{String} y \texttt{Real}. La restricción
\texttt{len("Serial Number") == 10} exige que el valor del número de serie
tenga exactamente diez caracteres. Las demás restricciones establecen que el
peso sea mayor que 30 y que el color no pueda tomar el valor
\texttt{'Pink'}.

\begin{lstlisting}[
    language={},
    caption={Modelo UVL completo con características del \textit{Type level} y \textit{string constraints}.},
    label={lst:modelo-uvl-tipos}
]
include
    Type.string-constraints

features
    Smartwatch
        mandatory
            String Color
            Real Weight
            String "Serial Number"

constraints
    len("Serial Number") == 10
    Weight > 30.0
    Color != 'Pink'
\end{lstlisting}

Los tres listados anteriores se han validado mediante \texttt{UVLReader} de
Flamapy antes de incorporarlos al documento. Esta comprobación confirma que \texttt{UVLReader} puede leer los archivos
presentados, pero no demuestra por sí sola que los modelos sean
satisfacibles;
la evaluación de esa propiedad se aborda de forma independiente en el
Capítulo~\ref{chap:resultados}.


\section{Ecosistema de herramientas y aplicación en variabilidad software}
\label{sec:ecosistema_herramientas}

A continuación se presentan las herramientas empleadas para gestionar y analizar modelos de características más allá de su simple representación, las cuales proporcionan capacidades para la creación, gestión y análisis automático de modelos, integrando diferentes técnicas de razonamiento y facilitando su adaptación a nuevos escenarios. Dentro de este ecosistema destaca Flamapy como \textit{framework} orientado al análisis automático de modelos de características mediante una arquitectura extensible. Asimismo, destaca UVLHub como plataforma destinada a la gestión y compartición de modelos definidos en UVL.

\subsection{Flamapy como \textit{framework} de análisis de variabilidad}
\label{subsec:flamapy_como_framework}

Flamapy es un \textit{framework} desarrollado en Python orientado al análisis automático de modelos de características dentro del ámbito de la ingeniería de líneas de productos software \cite{flama2023,flamapy}. Su objetivo principal es proporcionar un entorno común para trabajar con modelos de variabilidad, facilitando tanto su procesamiento como la aplicación de diferentes técnicas de análisis automático. Para ello, Flamapy puede entenderse desde dos perspectivas complementarias: como una herramienta que permite ejecutar análisis sobre modelos de características de forma directa, y como un \textit{framework} extensible que proporciona una arquitectura modular para incorporar nuevos modelos, operaciones y técnicas de análisis.

Desde la perspectiva de herramienta, Flamapy proporciona diferentes mecanismos para facilitar su utilización en distintos escenarios. Entre ellos se encuentra una interfaz de línea de comandos (CLI, \textit{Command-Line Interface}) que permite ejecutar análisis directamente sobre modelos de variabilidad, así como una fachada de Python (\textit{Python Facade}) que simplifica su integración dentro de otros desarrollos software. Además, ofrece interfaces adicionales como servicios REST (\textit{Representational State Transfer}) mediante Swagger y soporte para ejecución en entornos web mediante WebAssembly (WASM)~\cite{webassembly}, permitiendo utilizar sus capacidades de análisis desde diferentes aplicaciones y contextos.

Como \textit{framework}, Flamapy sigue una arquitectura modular basada en \textit{plugins}, los cuales permiten extender sus capacidades sin modificar el núcleo principal del sistema. Esta arquitectura hace mucho más sencilla la incorporación de nuevos lenguajes de modelado, metamodelos o técnicas de análisis mediante componentes independientes que pueden integrarse dentro del ecosistema existente. Gracias a este enfoque, es posible adaptar Flamapy a diferentes necesidades relacionadas con la ingeniería de variabilidad, incorporando soporte para nuevos tipos de modelos o herramientas de razonamiento.

Otro de los elementos principales que ofrece Flamapy son las operaciones (\textit{operations}), que representan las diferentes tareas de análisis que pueden ejecutarse sobre modelos de características. Cada operación define una funcionalidad concreta, como la comprobación de satisfacibilidad, la obtención de configuraciones válidas o el cálculo de propiedades del modelo, independientemente del formato en el que este haya sido definido. Esta separación permite reutilizar las mismas operaciones sobre diferentes tipos de modelos siempre que exista una representación compatible dentro del \textit{framework}.

Por otra parte, las transformaciones (\textit{transformations}) permiten convertir modelos entre diferentes representaciones dentro de Flamapy. Estas transformaciones actúan como mecanismo de conexión entre los distintos metamodelos soportados, permitiendo que los modelos puedan ser procesados por diferentes componentes del \textit{framework} (o incluso no del \textit{framework}) independientemente de su representación original. Por ejemplo, el componente \textit{UVLWriter}, que es el que se ha utilizado para descargar los modelos creados por el generador diseñado, permite transformar un modelo de características almacenado internamente en Flamapy a su representación textual en formato UVL. Otras de las transformaciones existentes, por nombrar algunas, son \textit{AFM Reader}, \textit{AFM Writer}, \textit{XML Reader}, \textit{JSON Writer}, \textit{SPLOT Writer} o \textit{FM to BDD}. Todo esto facilita la interoperabilidad entre lenguajes de variabilidad, metamodelos y herramientas de análisis.

\subsection{UVLHub como plataforma de gestión de modelos UVL}
\label{subsec:uvlhub_como_plataforma}

UVLHub es una plataforma orientada a la gestión de modelos de características definidos mediante el lenguaje UVL (\textit{Universal Variability Language}) dentro del ámbito de la ingeniería de líneas de productos software \cite{uvlhub2024,uvlhub}. Tiene como objetivo proporcionar un repositorio especializado que permita almacenar, organizar, compartir y analizar modelos de variabilidad, facilitando su reutilización en entornos de investigación y desarrollo software, así como la colaboración entre usuarios.

La razón principal de su existencia se debe a la complicada localización de los modelos de características empleados por los usuarios, ya que estos los almacenaban en diferentes repositorios, páginas personales o proyectos independientes, utilizando además diferentes formatos de representación, dificultando su búsqueda, reutilización y análisis. UVLHub sigue los principios de ciencia abierta (\textit{Open Science}), como accesibilidad, transparencia y colaboración, con la finalidad de compartir estos modelos de forma homogénea, haciendo más sencilla la colaboración entre usuarios.

Una de las características más relevantes de UVLHub es su capacidad de integración con otras herramientas y servicios relacionados con la gestión y análisis de modelos de variabilidad. Por una parte, la plataforma se integra con Zenodo, un repositorio abierto de investigación que permite almacenar los conjuntos de datos de forma permanente mediante la generación de identificadores DOI (\textit{Digital Object Identifier}), garantizando que los modelos publicados puedan seguir estando disponibles aunque la plataforma dejase de mantenerse en el futuro. Por otra parte, UVLHub incorpora las capacidades de análisis automático proporcionadas por Flamapy, permitiendo obtener información relevante sobre los modelos almacenados, como el número de características, la validez del modelo o el número de configuraciones posibles que pueden derivarse de ellos.

En la Figura~\ref{fig:arquitectura_uvlhub} se puede ver una vista general sobre la arquitectura del sistema.

\begin{figure}[H]
    \centering
    \includegraphics[width=0.8\textwidth]{arquitectura_uvlhub.png}
    \caption{Arquitectura general de UVLHub y sus integraciones principales. Extraído de \cite{egc2024_instalacion_uvlhub}.}
    \label{fig:arquitectura_uvlhub}
\end{figure}

Además de sus integraciones externas, UVLHub organiza sus funcionalidades (\textit{features}) en módulos orientados a cubrir las principales necesidades relacionadas con la gestión de modelos UVL. Algunas de las más relevantes son:

\begin{itemize}
    \item \textbf{dataset}: representa el núcleo principal de UVLHub, permitiendo gestionar el ciclo de vida de los conjuntos de modelos UVL. Esta funcionalidad permite realizar operaciones como la subida de datasets, su almacenamiento junto con los metadatos asociados, la descarga de los modelos y la resolución de identificadores DOI generados mediante la integración con Zenodo.

    \item \textbf{featuremodel}: permite gestionar la información asociada a los modelos de características almacenados dentro de un dataset. Esta funcionalidad define elementos como el modelo de características, sus metadatos y las métricas asociadas, estableciendo la relación entre los datasets y los archivos UVL almacenados.

    \item \textbf{explore}: proporciona mecanismos para explorar los datasets disponibles dentro de la plataforma mediante diferentes criterios de búsqueda, permitiendo localizar modelos de características publicados previamente.

    \item \textbf{flamapy}: integra las capacidades de análisis automático proporcionadas por Flamapy dentro de UVLHub. Esta integración permite realizar operaciones sobre los modelos UVL, como la validación del modelo o transformaciones hacia otros formatos soportados.

    \item \textbf{zenodo}: permite la comunicación con la interfaz de programación de aplicaciones (API, \textit{Application Programming Interface}) REST de Zenodo para realizar la publicación y almacenamiento permanente de los datasets, proporcionando identificadores DOI asociados a los recursos compartidos.
\end{itemize}

La Figura~\ref{fig:modulos_uvlhub} representa los módulos actuales de UVLHub en conjunto y cómo se complementan entre sí:

\begin{figure}[H]
    \centering
    \includegraphics[width=0.8\textwidth]{modulos_uvlhub.png}
    \caption{Estructura de módulos de UVLHub. Extraído de \cite{egc2024_instalacion_uvlhub}.}
    \label{fig:modulos_uvlhub}
\end{figure}

Estos módulos son originales de UVLHub. A ellos hay que añadirle también el módulo \textbf{generator}, del cual también se disponía antes de este trabajo. 

\subsection{Generación automática de modelos de características}
\label{sec:concepto4}

Los modelos de características permiten representar la
variabilidad existente dentro de una línea de productos software. Sin embargo,
a medida que aumenta el número de características, restricciones y posibles
configuraciones, la creación manual de estos modelos puede convertirse en una
tarea costosa y difícil de gestionar.

Por este motivo, la generación automática de modelos de características se ha
utilizado como mecanismo para producir modelos sintéticos destinados, entre
otros usos, a la evaluación y prueba de herramientas de análisis de variabilidad
\cite{segura2011metamorphic,galindo2023llm}.

Estas técnicas permiten obtener conjuntos de modelos con diferentes tamaños,
estructuras y grados de complejidad, facilitando la realización de experimentos
sobre distintos escenarios de evaluación.

\section{Trabajos relacionados}
\label{sec:relacionados}



Dentro del ámbito de la ingeniería de líneas de productos software existen
diferentes propuestas orientadas a la generación automática de modelos de
características. Estas aproximaciones se emplean principalmente para crear
modelos sintéticos destinados a evaluar herramientas de análisis de
variabilidad, realizar experimentos controlados o estudiar la complejidad de
técnicas de razonamiento.

La comparación se organiza según cuatro criterios: las construcciones de UVL
que puede generar cada propuesta, el mayor tamaño de modelo documentado, la
existencia de una garantía de satisfacibilidad y la validez sintáctica de la
salida. Los dos últimos criterios permiten distinguir entre producir un modelo
con una determinada estructura, garantizar que posee una configuración válida y
comprobar que el artefacto resultante puede ser interpretado por un parser UVL.
Además de Sundermann et al.~\cite{sundermann2024generating}, se consideran
S.P.L.O.T., FeatureIDE, BeTTy y ETHOM, por representar enfoques con objetivos
distintos dentro de la generación y evaluación de modelos de características.

\subsection{Generating Feature Models with UVL's Full Expressiveness}

El trabajo de Sundermann et al.~\cite{sundermann2024generating} constituye la
base conceptual del presente TFG, tal como se ha expuesto en la
Sección~\ref{sec:punto_partida}. Los autores presentan un generador de modelos
UVL que cubre todos los niveles de expresividad del lenguaje, incluyendo
atributos, expresiones aritméticas, tipos, cadenas y funciones agregadas.

La propuesta permite configurar el número de modelos y distintos parámetros
estructurales de la generación. También incorpora de forma opcional una
garantía de satisfacibilidad módulo teorías (SMT,
\textit{Satisfiability Modulo Theories}), aplicable a las construcciones que
soporta. El artículo muestra configuraciones de ejemplo con entre 800 y 1000
características, pero no establece un tamaño máximo evaluado ni informa de una
validación de la salida mediante un parser UVL externo.

Por tanto, este TFG no pretende superar a dicho trabajo en expresividad del
lenguaje. Su aportación se sitúa en la implementación e integración de un
componente compatible con Flamapy dentro de UVLHub, así como en la validación
de los modelos generados mediante \texttt{UVLReader}. La garantía de
satisfacibilidad implementada en este trabajo se limita actualmente a modelos
booleanos, por lo que constituye una limitación frente al enfoque SMT de
Sundermann et al.

\subsection{S.P.L.O.T.}

S.P.L.O.T. (\textit{Software Product Lines Online Tools}) es una plataforma
web orientada al desarrollo y análisis de líneas de productos software que
incluye un generador de modelos de características sintéticos
\cite{splot}. Sus modelos se basan en estructuras jerárquicas y restricciones
lógicas del nivel booleano.

En consecuencia, no contempla atributos tipados, expresiones aritméticas,
restricciones sobre cadenas ni otras construcciones avanzadas de UVL. La
documentación consultada no establece un tamaño máximo de modelo ni una opción
para garantizar la satisfacibilidad de todas las salidas generadas. Además, su
formato de salida no es UVL, por lo que no se documenta una validación sintáctica
mediante un parser de dicho lenguaje.

\subsection{FeatureIDE}

FeatureIDE es un \textit{framework} extensible para el desarrollo software basado en
características, que proporciona herramientas para crear, editar y analizar
modelos de características \cite{featureide}. También incorpora un generador
de modelos sintéticos integrado en su propio ecosistema.

En el contexto de los generadores comparados por Sundermann et
al.~\cite{sundermann2024generating}, la generación de FeatureIDE se limita
principalmente a construcciones booleanas. Por ello, no proporciona generación
de atributos ni de restricciones propias de los niveles avanzados de UVL. No se
ha identificado un tamaño máximo documentado, una garantía de
satisfacibilidad de la salida ni una validación sintáctica mediante un parser
UVL, puesto que los modelos se generan en el formato propio de la herramienta.

\subsection{BeTTy}

BeTTy (\textit{Benchmarking and Testing on the Automated Analysis of Feature
Models}) es un \textit{framework} orientado al \textit{benchmarking} y a la evaluación
de herramientas de análisis automático de modelos de características
\cite{betty}. Permite configurar parámetros de la generación para crear
escenarios experimentales con diferente complejidad estructural y distinto
número de restricciones.

BeTTy permite generar tanto modelos de características tradicionales como
modelos con atributos. Además, documenta una opción para solicitar modelos
satisfacibles en su generador de modelos clásico. Estas capacidades no
equivalen a generar archivos UVL que incorporen todos los niveles del
lenguaje: la documentación consultada no acredita la producción de
características tipadas y restricciones sobre cadenas conforme a UVL, ni una
validación de la salida mediante un parser de dicho lenguaje. Su aportación
se centra en la configuración de modelos destinados a la evaluación de
herramientas de análisis.

\subsection{ETHOM}

ETHOM (\textit{Evolutionary algoriTHm for Optimized feature Models}) es una
propuesta basada en algoritmos evolutivos para generar modelos de
características con propiedades que resultan exigentes para las herramientas de
análisis \cite{ethom}. A diferencia de los generadores aleatorios, formula la
obtención del modelo como un problema de optimización y persigue, por ejemplo,
maximizar el tiempo de análisis o el consumo de memoria. Su prototipo se desarrolló para integrarse en el entorno de BeTTy, por lo
que aquí se estudia como un enfoque de generación evolutiva dentro de ese
ecosistema \cite{ethom}.

ETHOM se incluye en esta revisión porque comparte la finalidad general de crear
casos de evaluación para herramientas de variabilidad. Sin embargo, su objetivo
no es generar modelos UVL con expresividad avanzada, sino modelos de
características tradicionales difíciles de analizar. Por ello, los criterios de
construcciones UVL y validez sintáctica UVL no son aplicables directamente; la
propuesta tampoco documenta un límite de tamaño ni una garantía general de
satisfacibilidad equivalente a la considerada en este trabajo.

\subsection{Comparativa de trabajos relacionados}
\label{subsec:comparativa_relacionados}

La Tabla~\ref{tab:comparativa_generadores} reúne las capacidades documentadas de cada propuesta y los resultados comprobados para el generador de este TFG.

\begin{table}[H]
\centering
\small
\setlength{\tabcolsep}{4pt}
\renewcommand{\arraystretch}{1.18}
\begin{tabularx}{\textwidth}{@{}
  >{\raggedright\arraybackslash}p{2.4cm}
  >{\raggedright\arraybackslash}p{3.6cm}
  >{\raggedright\arraybackslash}X
@{}}
\toprule
\textbf{Trabajo} & \textbf{Construcciones UVL} &
\textbf{Tamaño, satisfacibilidad y salida} \\
\midrule
S.P.L.O.T. \cite{splot} & Booleanas. &
Tamaño y satisfacibilidad no documentados. Formato propio, sin validación
UVL documentada. \\

FeatureIDE \cite{featureide} & Booleanas en su generador. &
Tamaño y satisfacibilidad no documentados. Formato propio, sin validación
UVL documentada. \\

BeTTy \cite{betty} &
Modelos booleanos y modelos con atributos; cobertura completa de los
niveles UVL no documentada. &
Tamaño máximo no establecido en la fuente comparada. Opción documentada
para solicitar modelos satisfacibles. Formatos distintos de UVL, sin
validación UVL documentada. \\

ETHOM \cite{ethom} & Modelos tradicionales; no orientado a UVL. &
Tamaño y satisfacibilidad no documentados. La validez sintáctica UVL no es
aplicable porque no genera salida UVL. \\

Sundermann et al. \cite{sundermann2024generating} & Todos los niveles de UVL. &
Ejemplo de 800--1000 características, sin máximo declarado. Satisfacibilidad
opcional mediante SMT. Genera UVL, sin validación externa mediante parser
documentada. \\

\midrule

\textbf{Este TFG} & Niveles booleano, aritmético y de tipos, con las construcciones y niveles menores implementados en el generador. &
100 características por modelo evaluadas. Satisfacibilidad parcial: solo
modelos booleanos mediante PySAT. Los 30 modelos evaluados fueron aceptados
por \texttt{UVLReader}. \\
\bottomrule
\end{tabularx}
\caption{Comparación de los trabajos relacionados según la expresividad,
el tamaño documentado, la satisfacibilidad y la validez sintáctica de la
salida. ``No documentado'' indica que la fuente consultada no proporciona
esa evidencia.}
\label{tab:comparativa_generadores}
\end{table}

Esta tabla muestra que S.P.L.O.T. y FeatureIDE se orientan principalmente
a construcciones booleanas, mientras que BeTTy también permite generar
modelos con atributos, aunque la tabla no documenta cobertura completa
de los niveles de UVL. ETHOM ofrece un enfoque diferente, pues
optimiza la dificultad de los modelos generados, aunque no persigue producir
artefactos UVL ni ampliar la expresividad del lenguaje.

La propuesta de Sundermann et al. es la referencia más próxima, al cubrir todos
los niveles de UVL y ofrecer una garantía opcional de satisfacibilidad. Frente a
ella, este TFG aporta la integración del generador en Flamapy y UVLHub y una
comprobación explícita de la validez sintáctica de los modelos generados con
\texttt{UVLReader}. No obstante, no se afirma una ventaja en tamaño máximo
evaluado ni en garantía de satisfacibilidad, ya que esta última se limita a la
configuración booleana actual.


El presente capítulo ha establecido los fundamentos teóricos y tecnológicos
necesarios para comprender el contexto del trabajo, así como las principales
soluciones existentes relacionadas con la generación automática de modelos de
características. A partir de este análisis se define el marco sobre el que se
desarrolla la propuesta presentada en este TFG. El siguiente capítulo describe
la planificación seguida durante el proyecto, incluyendo la metodología,
las fases de desarrollo, los recursos empleados y la gestión de riesgos.